\section{Introduction}
\label{sec:intro}
Stream processing systems process an infinite amount of data in a continuous manner.  They are essential components in
many modern applications. For example, Internet of Things (IoT) applications that process continuous streams of data
generated by large numbers of devices~\cite{sasaki2021survey}. Other examples include e-commerce websites~\cite{tucker2008nexmark}, social
networks~\cite{noghabi2017samza, toshniwal2014storm}, and network monitoring systems~\cite{bhatia2021mstream}, which process large volumes of data generated at high
speed.

A query in a stream processing system is represented by a directed acyclic graph (DAG) of {\em operators}, for example, Map, Filter, and Join operators. Data flows through the DAG, and each operator can maintain some states during query
execution. State-of-the-art stream processing systems partition the DAG over multiple machines, such that operators can
be executed in parallel. Examples of such distributed stream processing systems (DSP) include platforms such as Flink \cite{carbone2015apache},
Storm \cite{toshniwal2014storm}, and Spark Streaming \cite{zaharia2013discretized}. These systems are able to scale up to handle large and high-velocity streams while ensuring correctness and fault tolerance~\cite{lin2016streamscope}. Most of them run on a cloud platform to leverage the latter's resources and
reliability. There are data streaming services, which are built on top of DSP, offered by major cloud
providers, for example, AWS's Kinesis and Azure's Stream Analytics.

Securing DPS against untrusted cloud providers requires protecting both the data {\em and the queries}. In general, the data stream may contain sensitive information, such as physiological readings, locations, or network activities. 
%
However, in many applications, the query executed over the stream can even be more sensitive and critical when they represent proprietary algorithms for data processing. One example is Bargain Index \cite{spade}, which finds bargains in the stock trading market. The query ingests a trading stream and a quote stream (sell orders), then computes bargain indexes. This query enables financial institutions to make trading decisions, thus it is an important business asset. An attacker (e.g. a competitor) knowing the query can steal and even optimize it to make faster, better trading decisions to gain profits at the expense of the victim. Another example is Network Spike Detection \cite{SAQL}, which tracks network usage and raises alerts when it exceeds a predefined threshold, indicating potential data exfiltration. An attacker knowing the query, especially the predefined threshold, can adapt its behavior to avoid detection. 

% For example, consider a network monitoring application that runs a query for detecting anomalous events from a stream of
% network activities. 
% %
% An attacker learning the content of the query or tampering with the query execution can evade the
% detection. 
% %
% As another example, consider an e-commerce application that runs a query implementing an extract-transform-load (ETL) pipeline to extract features from user transactions for a downstream market prediction
% model~\cite{de2024hopsworks}. 
% %
% An attacker learning this pipeline can undermine the company's business advantage and even reveal such strategic logical operations to competitors.
% %
% Furthermore, in existing streaming healthcare biosensor data analytic applications~\cite{chen2004stream}, leaking which biosensor queries are being processed can reveal certain health issues and safety concerns of patients.

Existing works on protecting outsourced database systems, of which DPS is an instance, often leverage homomorphic encryption, secure multiparty computation, 
trusted hardware, or a combination thereof \cite{schuster2015vc3, popa2011cryptdb,reijsbergen2023tap,dinh2015m2r,zhang2015integridb}.
%
Given the stringent latency requirement
in stream processing, state-of-the-art secure DPS systems typically use trusted hardware  \cite{DBLP:journals/corr/abs-1805-01752,
234964,schuster2015vc3, dauterman2022waldo}. These systems run the stream engine inside trusted execution environments (TEEs), e.g., Intel SGX enclaves. 
%
The data is encrypted and is only processed in plaintext inside the TEEs.

In this study, we show that DSP systems that run vanilla stream engine inside the trusted hardware are vulnerable to
timing channel attacks. Specifically, we consider systems similar to Flink \cite{carbone2015apache}, which run the query DAG inside the 
Intel SGX enclaves. Our attack leverages two key insights. 
%
First, because of the distributed design, it is possible to
isolate individual operators from the query DAG and to run them in separate enclaves. 
%
By monitoring the enclave
execution, the attacker can extract fine-grained timing information of each operator. Second, the attacker can run their
own enclaves that execute {\em similar} operators with {\em similar} data to those of the target query. This is because
the data schema can be known by the attacker, for example, the attacker being a valid user of the application.
Furthermore, most DSP software is modular and open source, allowing the attacker to create enclaves for operators of
their choice and execute them with their own synthetic data.

\begin{figure}
    \centering
    \includegraphics[width=0.49\textwidth]{Figures/attack_flow.pdf}
    \caption{Overview of the attack, with an offline phase and an online phase.}
    \label{fig:attack-flow}
\end{figure}

Figure~\ref{fig:attack-flow} illustrates how the attacker exploits the two insights above to leak the content of the
query DAG. The attack occurs in two phases. In the offline phase, the attacker profiles the timing characteristics of
individual operators. 
%
In particular, for each operator, it creates an enclave and executes it with synthetic data while
recording the execution time. The result is a large, labeled dataset of execution times. At the end of this phase, the
attacker trains a machine learning model that predicts the operator based on the execution times. In the online phase,
the attacker isolates individual operators from the target query, such that each operator is executed in a separate
enclave. This is possible because the attacker has control over resources that the DSP system assigns to the query. It
monitors each enclave's execution and uses the model in the offline phase to learn the content and the parameters of the
operator.

We implement our attack against a customized DSP execution engine. We use popular streaming benchmarks to generate
synthetic data and queries for the offline phase. We test the trained models against unknown queries from the same
benchmarks. The results demonstrate that even with simple models, the operators can be leaked with high accuracy, of up
to $89\%$. The models can also recover the parameters inside each operator with high accuracy. The main cost
of the attack is the offline phase, but it can be significantly reduced because this phase is highly parallelizable. Our
results suggest that timing channels are a viable attack vector against DSP systems. We discuss system-level,
lightweight mitigations that undermine the offline phase by batch scheduling or operator merging.

In summary, we make the following contributions:
\begin{itemize}
\item We propose a side channel attack on distributed stream processing systems that runs a vanilla execution engine
inside Intel SGX. The attack is based on profiling the execution time of individual query operators.
\item We demonstrate that the attack can leak query operators with high accuracy, using popular streaming benchmarks.
The artifact is published at \url{https://github.com/PhamHung2020/Simple-SGX-based-DSP-Engine}.
\item We discuss lightweight mitigation approaches to harden against attacks by avoiding the collection of precise timing information, thereby making it more difficult to leak query operators.
\end{itemize}

The remainder of this paper is organized as follows. Section \ref{sec:background} provides the background on distributed
stream processing systems and Intel SGX. Section~\ref{sec:model} follows with the system and threat model.
Section~\ref{sec:methodology} details the timing side-channel attack. Section~\ref{sec:experiments} evaluates the
attack's success rate. Section~\ref{sec:discussion} discusses mitigation approaches and other limitations of the attack.
Section~\ref{sec:related} discusses the related work before Section~\ref{sec:conclusions} concludes.






